\section{Implementação do Sistema de Snapshot e Auditoria}
\label{sec:snapshot-auditoria}
\subsection{Entidade Snapshot e Modelo de Persistência}
\label{sec:entidade-snapshot}

A entidade \texttt{SpecialistSnapshot} foi criada para preservar o estado 
do especialista antes da aplicação de alterações aprovadas no fluxo de 
mudança de categoria. O nome reflete o escopo: o snapshot cobre o 
especialista, e é acionado exclusivamente pelo tipo \texttt{CHANGE\_CATEGORY} 
— recurso e atualização de pontuação não geram snapshot, pelas razões 
discutidas na Seção~\ref{sec:efeitos-por-tipo}. O 
Quadro~\ref{qua:modelo-snapshot} descreve sua estrutura.

\begin{quadro}[htbp]
\centering
\caption[Modelo da entidade SpecialistSnapshot]{Modelo da entidade \texttt{SpecialistSnapshot}.}
\label{qua:modelo-snapshot}
\small
\setlength{\tabcolsep}{4pt}
\begin{tabular}{p{2.8cm}p{2cm}p{3.4cm}p{4.4cm}}
\toprule
\textbf{Coluna} & \textbf{Tipo} & \textbf{Restrição} & \textbf{Observação} \\
\midrule
\texttt{id} & Inteiro & Chave primária & Identificador do snapshot \\
\texttt{specialist\_id} & Inteiro & Obrigatório, indexado & Especialista cujo estado foi capturado; sem FK \\
\texttt{snapshot} & JSONB & Obrigatório & Estado anterior serializado \\
\texttt{authorizer\_id} & Inteiro & Obrigatório & Usuário que autorizou a alteração; sem FK \\
\texttt{created\_at} & DataHora & Padrão \texttt{agora()} & Momento da captura \\
\bottomrule
\end{tabular}
\fonte{Elaboração própria (2026).}
\end{quadro}

Duas decisões de modelagem merecem destaque. A primeira é a 
\textbf{ausência de chave estrangeira} em \texttt{specialist\_id} e 
\texttt{authorizer\_id}. Trata-se de uma escolha deliberada para uma 
entidade de auditoria: o snapshot é um registro histórico, e sua 
integridade não deve depender da existência atual da entidade referenciada. 
Se o especialista ou o autor da decisão for removido do sistema, o 
snapshot permanece consultável, preservando a memória do que ocorreu. A 
contrapartida é a possibilidade de referências órfãs e a ausência de 
integridade referencial garantida pelo banco — custo aceito em troca da 
preservação histórica. A alternativa — FK com \texttt{ON DELETE SET NULL} 
— destruiria justamente a informação que o snapshot existe para preservar.

A segunda é a \textbf{ausência de vínculo explícito com a solicitação} 
que originou o snapshot. A entidade referencia o especialista, não a 
solicitação nem a revisão. Essa decisão tem consequência direta sobre a 
reversão (Seção~\ref{sec:reversao}): a operação de reversão busca o 
snapshot mais recente do especialista, e não o snapshot "daquela 
solicitação". Para o fluxo atual — em que se reverte a última alteração 
aprovada —, a modelagem é suficiente. Para reversão de uma alteração 
específica no meio do histórico, seria necessário vincular o snapshot 
à solicitação ou à revisão correspondente, o que está registrado como 
limitação em Seção~\ref{sec:rastreabilidade}. 

O campo \texttt{snapshot} armazena o estado anterior em JSONB. A escolha 
do formato, em vez de colunas relacionais, decorre de três fatores. O 
primeiro é a \textbf{coesão do estado}: os campos capturados descrevem um 
único estado lógico do especialista e fazem sentido juntos; armazená-los 
em uma coluna única preserva essa coesão e evita reconstituição por 
junção. O segundo é a \textbf{evolução do esquema}: se o conjunto de 
campos a capturar for ampliado, um documento JSONB acomoda novos campos 
sem migration; colunas adicionais exigiriam alteração de esquema e 
migração de dados históricos. O terceiro é a \textbf{legibilidade}: 
consultas diretas ao documento permitem visualizar o estado anterior 
sem conhecer o esquema das tabelas de origem.

As contrapartidas também precisam ser declaradas. O JSONB não impõe 
integridade referencial interna — o documento é um retrato, não uma 
referência viva, e não há como garantir que os campos correspondam ao 
esquema atual do especialista depois de migrations futuras. O formato 
não é validado pelo banco: a estrutura do documento depende da aplicação 
que o grava. Essas limitações são aceitáveis no escopo do módulo, porque 
o snapshot é um registro histórico e **deve** preservar o estado como 
era no momento da captura, e não acompanhar a evolução do esquema. A 
alternativa — colunas relacionais com FK para as tabelas de origem — 
traria integridade, mas produziria snapshots que se corrompem quando a 
tabela de origem muda, anulando o propósito de preservar o estado anterior.

A captura é feita antes da aplicação dos efeitos e dentro da mesma 
unidade de trabalho da aprovação (Seção~\ref{sec:efeitos-por-tipo}): o 
snapshot, a mudança de categoria e o registro da revisão são confirmados 
em conjunto. Se qualquer etapa falhar, o snapshot é descartado junto, 
garantindo que ele exista se e somente se a alteração correspondente foi 
efetivada. Isso atende ao RNF01 e sustenta o RNF02: sem atomicidade, um 
snapshot poderia registrar um estado anterior que nunca foi alterado, ou 
uma alteração sem snapshot correspondente.

O campo \texttt{authorizer\_id} registra o usuário responsável pela 
autorização da alteração, separadamente do autor da revisão. Na prática, 
os dois coincidem no fluxo atual — o avaliador que aprova a revisão é o 
mesmo que autoriza a alteração —, mas manter o campo explícito na 
entidade permite que, no futuro, a autorização seja delegada ou 
automatizada sem alteração do modelo. A separação também garante que a 
informação esteja presente na própria entidade de auditoria, e não apenas 
derivável por junção com a revisão.

\subsection{Reversão de Alterações}
\label{sec:reversao}

A reversão permite desfazer uma alteração de categoria aprovada, 
restaurando o estado do especialista registrado no snapshot mais recente. 
A operação é acionada por rota administrativa e está sujeita a controle 
de acesso, como as demais operações de decisão (RNF03).

A reversão não é uma operação destrutiva: o snapshot não é apagado após 
o uso, e a reversão é registrada como evento próprio na linha do tempo 
(Seção~\ref{sec:timeline}). O histórico completo é preservado — a 
solicitação original, a revisão que aprovou, o snapshot capturado, a 
reversão e o estado resultante. Essa é a diferença entre reversão e 
exclusão de histórico: a primeira adiciona um evento, a segunda remove 
o passado. A escolha pela primeira é o que sustenta a auditoria.

O Algoritmo~\ref{alg:reversao} descreve o fluxo em pseudocódigo, com 
base na recuperação do snapshot mais recente do especialista.

\begin{algorithm}[htbp]
\caption{Fluxo de reversão (pseudocódigo).}
\label{alg:reversao}
\begin{verbatim}
FUNÇÃO reverter(especialista_id, autorizador):
    especialista = buscar Specialist por id
    SE não existe:
        erro "especialista não encontrado"

    snapshot = buscar SpecialistSnapshot
               filtrando specialist_id = especialista.id
               ordenando por created_at DESC
               limitando a 1 resultado
    SE não existe:
        erro "nenhum snapshot encontrado"

    estado_anterior = desserializar(snapshot.snapshot)

    EM UMA TRANSAÇÃO:
        restaurar(especialista, estado_anterior)
        registrar_evento(especialista, REVERTIDA,
                         autorizador, data)

    RETORNAR confirmação
\end{verbatim}
\end{algorithm}

Quatro decisões delimitam o escopo da reversão. A primeira é que ela 
restaura o \textbf{snapshot mais recente}, não o snapshot de uma 
solicitação específica. Isso decorre da modelagem da entidade 
(Seção~\ref{sec:entidade-snapshot}), que referencia o especialista sem 
vincular-se à solicitação. A consequência é que a reversão desfaz a 
última alteração aprovada, e não uma alteração qualquer no histórico. 
Para o caso de uso original — corrigir uma mudança de categoria recente 
que se mostrou inadequada —, a semântica é adequada. Para reverter uma 
alteração específica no meio do histórico, seria necessário alterar o 
modelo para vincular o snapshot à solicitação. A limitação está 
registrada em Seção~\ref{sec:rastreabilidade}.

A segunda é que a reversão restaura apenas os campos capturados no 
snapshot, sem desfazer alterações feitas ao especialista por outros 
fluxos no intervalo. Se, após a aprovação, o especialista atualizou 
dados pessoais ou pontuação por outro caminho, esses dados não são 
afetados. A escolha é conservadora e evita que uma reversão de categoria 
apague, por tabela, alterações não relacionadas.

A terceira é que a reversão \textbf{não altera o estado da revisão} que 
aprovou a solicitação. A revisão permanece registrada como aprovada, e 
a reversão é um evento subsequente — o que preserva o histórico de 
decisões.

A quarta é que a operação adiciona um novo evento à linha do tempo, mas 
não gera um novo snapshot. Isso poderia ser questionado — se a reversão 
também é uma alteração, por que não capturar o estado anterior a ela? A 
resposta é que reverter uma reversão seria equivalente a reaplicar a 
alteração original, e isso é feito criando-se uma nova solicitação, não 
reusando o snapshot antigo. A decisão mantém o modelo simples: um 
snapshot por alteração aprovada, um evento de reversão por reversão 
executada.

\subsection{Linha do Tempo de Eventos}
\label{sec:timeline}

A linha do tempo é a apresentação cronológica dos eventos associados a 
um especialista: criação da solicitação, submissão para análise, decisão 
da revisão, captura de snapshot e, quando aplicável, reversão. A 
implementação \textbf{não mantém tabela própria de eventos}: a linha do 
tempo é derivada das entidades que já registram cada evento com autor e 
data — \texttt{Solicitation} (\texttt{created\_at}), 
\texttt{SolicitationReview} (autor, veredicto, comentário, data) e 
\texttt{SpecialistSnapshot} (autorização, captura, data). A escolha por 
derivação, em vez de uma tabela de log genérica, evita duplicação de 
informação e garante que a linha do tempo nunca divirja das entidades de 
origem: se a revisão registrou uma decisão, ela aparece na linha do 
tempo; se o snapshot foi capturado, ele também aparece.

O preço da derivação é que a consulta precisa agregar três fontes e 
ordená-las por data. Isso é feito no serviço de consulta da linha do 
tempo, que retorna a lista consolidada para a interface. A alternativa 
— tabela de eventos dedicada, alimentada por gatilhos — teria a vantagem 
de uma consulta única, mas exigiria sincronização com as entidades de 
origem e abriria margem para inconsistência. Para o volume atual, a 
derivação é a solução mais simples e mais robusta.

A linha do tempo cumpre duas funções. Para o solicitante, mostra o 
andamento da solicitação e a justificativa de cada decisão. Para o 
avaliador e para auditoria, permite reconstituir o processo completo, 
com autores, datas e estados intermediários. Em ambos os casos, a 
informação apresentada é a mesma que sustenta a rastreabilidade — a 
interface apenas a torna legível.

\subsection{Rastreabilidade e Limitações}
\label{sec:rastreabilidade}

O conjunto formado por solicitação, revisão, snapshot e linha do tempo 
atende ao requisito de rastreabilidade (RNF02) em três dimensões. A 
dimensão \textbf{temporal} é coberta pelos carimbos automáticos das 
entidades. A dimensão \textbf{autoral} é coberta pelo registro do 
solicitante na solicitação, do avaliador na revisão e do autorizador no 
snapshot. A dimensão \textbf{semântica} — o que mudou e por quê — é 
coberta pelo snapshot (estado anterior) e pela revisão (justificativa). 
As três juntas permitem reconstituir, para qualquer solicitação, quem 
pediu, quem decidiu, quando, com qual justificativa e de qual estado 
para qual estado.

O mecanismo tem limitações que precisam ser declaradas. A primeira é a 
já discutida \textbf{reversão pelo snapshot mais recente}: como o 
snapshot não referencia a solicitação, a reversão desfaz a última 
alteração aprovada, e não uma alteração específica. A segunda é a 
\textbf{ausência de exclusão lógica}: o \texttt{cascade="all, 
delete-orphan"} na relação entre solicitação e revisões implica que a 
remoção física de uma solicitação apaga o histórico associado. O módulo 
não expõe rota de exclusão, de modo que o cascade atua apenas como 
salvaguarda de integridade, mas a ausência de uma política de retenção 
é uma lacuna que pode ser endereçada em trabalho futuro. A terceira é 
a \textbf{ausência de vínculo do snapshot com a solicitação}: o 
snapshot referencia o especialista, não a solicitação, o que simplifica 
a modelagem e atende ao caso de uso atual, mas impede auditoria 
retroativa que relacione snapshot e revisão diretamente — a associação 
existe apenas por proximidade temporal e autoral. A quarta é que 
\textbf{alterações fora do fluxo formal não geram snapshot}: 
intervenções diretas no banco, ainda que possíveis para a equipe de 
suporte, não geram revisão nem snapshot, apenas atualizam 
\texttt{last\_update}. A política de acesso ao banco em produção está 
fora do escopo do módulo, mas o efeito sobre a rastreabilidade precisa 
ser reconhecido.

Essas limitações não invalidam o mecanismo; delimitam o que ele garante. 
O sistema passou de um fluxo sem registro formal (Seção~\ref{sec:fluxo-antigo}) 
para um fluxo em que toda operação realizada pela aplicação é registrada 
com autor, data, justificativa e estado anterior. O que está fora do 
fluxo formal — intervenções diretas no banco — permanece fora do escopo 
da garantia, e isso é uma decisão de escopo, não uma falha do mecanismo.